\documentclass[preprint,12pt]{elsarticle}
\usepackage{booktabs}
\usepackage{amsmath,amssymb}
\usepackage{graphicx}
\usepackage{hyperref}

\journal{Neural Computing and Applications}

\begin{document}
\begin{frontmatter}
\title{Physics-Guided Satellite Data Assimilation with Ablation-Oriented Evaluation}
\author{First Author}
\address{School/Institute, City, Country}
\begin{abstract}
This draft reports an ablation-oriented neural assimilation pipeline.
In this version, V1 is treated as Ours, while Full is considered an additional ablation variant.
\end{abstract}
\begin{keyword}
Satellite data assimilation \sep Neural data-driven model \sep Ablation study
\end{keyword}
\end{frontmatter}

\section{Results}
\begin{table}[htbp]
\centering
\caption{Main results and ablation comparison.}
\begin{tabular}{llccccc}
\toprule
Method & Type & RMSE(K) & MAE(K) & Bias(K) & Corr & Improve. \\
\midrule
Background (ERA5) & bkg & 1.8230 & 1.2177 & +0.0020 & 0.99624 & +0.00\% \
OI/1DVar (B2) & oi & 1.8836 & 1.2798 & -0.0014 & nan & -3.32\% \
\textbf{Ours (V1)} & ours & 1.1696 & 0.8054 & +0.0229 & 0.99858 & +35.84\% \
Full Variant (Ablation) & ablation & 1.2674 & 0.8650 & +0.0596 & 0.99832 & +30.48\% \
w/o MaskConv (V2) & ablation & 1.2150 & 0.8337 & +0.0366 & 0.99849 & +33.35\% \
w/o GatedFusion (V3) & ablation & 1.3734 & 0.9445 & +0.0830 & 0.99736 & +24.66\% \
w/o SpectralStem (V4) & ablation & 1.4265 & 0.9690 & +0.0828 & 0.99732 & +21.75\% \
VanillaUNet (B3) & compare & 1.4491 & 1.0093 & +0.1664 & 0.99684 & +20.51\% \
FuXi-DA (B4) & compare & 1.4095 & 0.9995 & +0.1043 & 0.99654 & +22.68\% \
AttentionUNet (B5) & compare & 1.5092 & 1.0508 & +0.1778 & 0.99666 & +17.21\% \
PixelMLP (B6) & compare & 1.6959 & 1.1654 & +0.1388 & 0.99628 & +6.97\% \
ResUNet (B7) & compare & 1.5407 & 1.0784 & +0.1783 & 0.99630 & +15.49\% \
\bottomrule
\end{tabular}
\end{table}

\begin{figure}[htbp]
\centering
\includegraphics[width=0.9\linewidth]{fig_combined_rmse_improve.png}
\caption{Combined RMSE and improvement comparison across all evaluated methods.}
\label{fig:combined_rmse_improve}
\end{figure}

\begin{figure}[htbp]
\centering
\includegraphics[width=0.62\linewidth]{vertical_rmse_comparison.png}
\caption{Vertical RMSE profiles across pressure levels. V1 is the reported Ours; Full is treated as an auxiliary ablation variant.}
\label{fig:vertical_rmse}
\end{figure}

\section{Ablation Narrative}
Based on current experiments, V1 is positioned as Ours and used as the principal reference in the ablation analysis.
The original Full setting is reinterpreted as an additional ablation branch to keep the experimental narrative consistent with observed ranking.

\section{Conclusion}
This manuscript draft can be extended with detailed method, dataset, and uncertainty analysis sections.
\end{document}